\documentclass[12pt]{letter}
\usepackage[utf8]{inputenc}
\usepackage{geometry}
\geometry{margin=1in}

\begin{document}

\begin{letter}{Editor-in-Chief\\
Applied Mathematics and Computation}

\opening{Dear Editor,}

We submit our manuscript entitled \textbf{``Spatiotemporal spectral derivatives with forward stepwise BIC for robust partial differential equation discovery''} by Fuchang Wang and Huirong Cao for consideration as a regular article in \textit{Applied Mathematics and Computation}.

\textbf{Problem.} Data-driven PDE discovery from noisy data faces two practical difficulties. First, high-order spatial derivatives such as $u_{xxxx}$ amplify measurement noise. Second, sparse regression methods like STLSQ require a manually chosen threshold, which varies with noise level and problem. Existing spectral derivative methods are spectrally accurate but assume periodic boundaries; on non-periodic data they suffer Gibbs ringing.

\textbf{What we do.} We propose three connected changes:

\begin{enumerate}
\item \textbf{Spatiotemporal spectral derivatives.} Instead of transforming space at each time step independently, we apply a single 2D Fourier transform over the joint space-time grid. For a traveling wave $u(x,t)=f(x-ct)$, the spectrum lies along the line $k_t=-ck_x$ (Proposition~2), whereas boundary Gibbs noise spreads isotropically. A 2D filter can therefore separate signal from boundary artifacts. This is the key difference from spatial-only spectral methods (Schaeffer and McCalla, 2017).

\item \textbf{Spectral weak form.} The regression is set up in weak form: both sides of the PDE are integrated against compactly supported Gaussian test functions. This integration averages out high-frequency noise in the spectrally computed derivatives, particularly for high-order terms. We give a spectral convergence error bound (Proposition~1).

\item \textbf{Normalized forward stepwise BIC.} Each candidate term is scaled to unit $\ell_2$ norm before regression. Forward stepwise selection with BIC then chooses the support automatically, with no threshold grid search. Under standard restricted-eigenvalue and minimum-coefficient conditions, the selected support is consistent (Proposition~3).
\end{enumerate}

\textbf{What we find.} On the MDBench benchmark (four 1D PDEs at five noise levels, plus three 2D/3D problems):

\begin{itemize}
\item On the non-periodic advection equation---where spatial-only spectral derivatives fail completely (0/1 terms recovered)---our method recovers both terms exactly at clean data.
\item On Burgers, all true terms are recovered at every noise level including SNR~10.
\item On KdV and Kuramoto--Sivashinsky, exact recovery is achieved at three and two noise levels, respectively.
\item On periodic problems, accuracy is comparable to Weak SINDy while runtime is 10--20 times lower.
\end{itemize}

We compare against three baselines: finite-difference SR3, spatial-only spectral derivatives with exponential filtering, and Weak SINDy. The control-variable comparison with spatial-only spectral methods isolates the effect of the spatiotemporal transform.

\textbf{Why AMC.} The paper combines a numerical algorithm (spatiotemporal spectral differentiation), error analysis (three propositions), and benchmark experiments on a public dataset. It does not require GPU acceleration and runs on a standard CPU. We believe this fits the scope of \textit{Applied Mathematics and Computation}.

This manuscript has not been published elsewhere and is not under consideration by another journal. All authors have approved the submission. We declare no competing interests. Code is available at \texttt{https://github.com/fzmath/spectral-pde-discovery}.

Thank you for your consideration.

\closing{Sincerely,}

\vspace{0.5cm}
\textbf{Fuchang Wang}\\
School of Science, Emergency Management University\\
No. 465 College Road, Sanhe, Hebei 065201, China\\
Email: wangfuchang@cidp.edu.cn\\
\textit{Corresponding author}

\vspace{0.3cm}
\textbf{Huirong Cao}\\
School of Science, Langfang Normal University\\
No. 100 Ai-min Road, Langfang, Hebei 065000, China\\
Email: huirongcao@126.com

\end{letter}
\end{document}
